\documentclass[11pt]{article}
\usepackage[margin=1in]{geometry}
\usepackage[T1]{fontenc}
\usepackage[utf8]{inputenc}
\usepackage{lmodern}
\usepackage{amsmath, amssymb, mathtools}
\usepackage{hyperref}
\usepackage{enumitem}
\usepackage{parskip}

\hypersetup{
  colorlinks=true,
  linkcolor=blue,
  urlcolor=blue
}

\title{Algorithm and Model Changes}
\author{magnetic-field-inspired-navigation}
\date{2026-05-22}

\begin{document}
\maketitle

This document records major algorithm and model decisions in the cleanroom
\texttt{magnetic-field-inspired-navigation} implementation.

\textbf{Rule for this repo}

\begin{itemize}[leftmargin=*]
  \item Any major change to the controller, dynamics model, obstacle model, sensing abstraction, or mathematical formulation must be recorded in this file in the same change that introduces it.
\end{itemize}

The purpose is to keep a durable map between:

\begin{itemize}[leftmargin=*]
  \item the legacy ROS implementation in \texttt{double\_integrator}
  \item the cleanroom Python implementation
  \item intentional simplifications or deviations
\end{itemize}

\section*{Entry 2026-05-22: Initial Cleanroom Double-Integrator Reference}

\textbf{Status}

\begin{itemize}[leftmargin=*]
  \item Implemented in \texttt{src/mfinav/double\_integrator.py}
  \item Scenario runner in \texttt{scripts/run\_reference.py}
\end{itemize}

\textbf{Legacy source anchors}

\begin{itemize}[leftmargin=*]
  \item Goal controller and relaxation: \texttt{nodes/ataka\_listener.py}
  \item Duplicate goal controller: \texttt{nodes/multi\_agent\_listener.py}
  \item Magnetic avoidance, \texttt{field == 2}: \texttt{nodes/collision\_avoidance\_okt.py}
  \item Double-integrator rollout: \texttt{nodes/system\_motion.py}
\end{itemize}

\subsection*{What Was Implemented}

The current cleanroom reference includes:

\begin{itemize}[leftmargin=*]
  \item 2D double-integrator dynamics
  \item 2D circular obstacle abstraction
  \item goal-relaxation gain logic inspired by the ROS implementation
  \item magnetic-field-inspired obstacle avoidance based on the ROS \texttt{field=2} branch
  \item additive total command
\end{itemize}

\begin{align}
a &= a_{\text{goal}} + a_{\text{obs}} \\
v_{k+1} &= v_k + a_k \Delta t \\
p_{k+1} &= p_k + v_{k+1}\Delta t
\end{align}

\subsection*{Goal Relaxation Implemented}

Let
\begin{itemize}[leftmargin=*]
  \item $p \in \mathbb{R}^2$ be robot position
  \item $v \in \mathbb{R}^2$ be robot velocity
  \item $g \in \mathbb{R}^2$ be goal position
  \item $r_g = g - p$ be the goal vector
  \item $r_o$ be the closest obstacle vector
  \item $d_g = \lVert r_g \rVert$
  \item $d_o = \lVert r_o \rVert$
\end{itemize}

The relaxation scale used in the cleanroom implementation is

\begin{equation}
k_{\text{relax}}
=
\left(1 - e^{-d_o/1.5}\right)
\cdot
w_{\angle}
\cdot
w_{\text{prog}}
\cdot
w_{\text{exit}}
\end{equation}

where

\begin{equation}
w_{\angle}
=
1 - \cos \theta
=
1 - \frac{r_g^\top r_o}{\lVert r_g \rVert \lVert r_o \rVert}
\end{equation}

\begin{equation}
w_{\text{prog}}
=
\begin{cases}
1, & d_g \le d_{g,0} \\
\exp\left(-\dfrac{d_g - d_{g,0}}{0.1}\right), & d_g > d_{g,0}
\end{cases}
\end{equation}

and

\begin{equation}
w_{\text{exit}}
=
\exp\left(-\dfrac{d_g - d_{g,\min}}{0.1}\right)
\end{equation}

with:

\begin{itemize}[leftmargin=*]
  \item $d_{g,0}$: initial distance to the goal
  \item $d_{g,\min}$: smallest distance to the goal reached so far
\end{itemize}

This follows the structure of the ROS code around the \texttt{goal\_relax == 1} block.

\subsection*{Goal Controller Implemented}

For the far field, the current cleanroom implementation uses a simplified
speed-and-turn controller inspired by the ROS geometric block:

\begin{equation}
a_{\text{goal}} = a_{\text{attr}} + a_{\text{turn}}
\end{equation}

\begin{equation}
a_{\text{attr}}
=
-\left(k_p\,k_{\text{relax}}\right)
\left(\lVert v \rVert - v_{\max}\right)\hat{s}
\end{equation}

where $\hat{s}$ is:

\begin{itemize}[leftmargin=*]
  \item the normalized velocity direction if $\lVert v \rVert > 0$
  \item otherwise the normalized goal direction
\end{itemize}

The turning term is

\begin{equation}
a_{\text{turn}}
=
k_{\omega}\,k_{\text{relax}}\,\lVert v \rVert\,\hat{n}\,(\hat{n}^\top \hat{g})
\end{equation}

where:

\begin{itemize}[leftmargin=*]
  \item $\hat{g} = r_g / \lVert r_g \rVert$
  \item $\hat{n}$ is the left-perpendicular of the current heading
\end{itemize}

Near the goal, the controller switches to a PD-like attractive form:

\begin{equation}
a_{\text{goal}}
=
k_{p,\text{goal}} r_g - k_d v
\end{equation}

\subsection*{Magnetic-Field-Inspired Obstacle Term Implemented}

The cleanroom port uses the \texttt{field == 2} structure from the old code.

Let
\begin{itemize}[leftmargin=*]
  \item $r_o$ be the closest vector from robot to obstacle surface
  \item $d = \lVert r_o \rVert$
  \item $\hat{l}_a = v / \lVert v \rVert$ be the normalized robot motion direction
  \item $r_{\text{from obs}} = -r_o$
  \item $\hat{r} = r_{\text{from obs}} / \lVert r_{\text{from obs}} \rVert$
\end{itemize}

The obstacle current direction is the tangential projection of the robot motion:

\begin{equation}
\hat{t}
=
\hat{l}_a - (\hat{l}_a^\top \hat{r})\hat{r}
\end{equation}

If this becomes degenerate, the implementation falls back to a perpendicular
direction.

The 2D magnetic field surrogate is modeled as

\begin{equation}
B \propto \frac{\lVert v \rVert}{d}\,\hat{t}_{\perp}
\end{equation}

and the magnetic avoidance term is

\begin{equation}
F_{\text{mag}}
=
k_m \hat{l}_{a,\perp}(\hat{l}_{a,\perp}^\top B)
\end{equation}

This is a 2D analogue of the nested cross-product structure in the ROS code:

\begin{align}
B &= [\hat{t}]_{\times} v / d \\
F &= k_m [\hat{l}_a]_{\times} B
\end{align}

\subsection*{Added Repulsive Correction Implemented}

The cleanroom implementation also keeps the old additional repulsive term near
the obstacle:

\begin{equation}
F_{\text{rep}}
=
-k_r \left(\frac{1}{d} - \frac{1}{d_b}\right)\frac{r_o}{d}
\quad \text{for } d \le d_b
\end{equation}

Then it is projected to remove the component along the motion direction:

\begin{equation}
F_{\text{rep}}^{\perp}
=
F_{\text{rep}} - (F_{\text{rep}}^\top \hat{l}_a)\hat{l}_a
\end{equation}

The final obstacle term is:

\begin{equation}
a_{\text{obs}} = F_{\text{mag}} + F_{\text{rep}}^{\perp}
\end{equation}

\subsection*{Major Simplifications Relative To The Old ROS Code}

These are intentional and important:

\begin{enumerate}[leftmargin=*]
  \item The current reference is strictly 2D. The ROS code contains both 2D and 3D branches.
  \item Obstacles are currently circles. The ROS stack computes obstacle geometry from point clouds, polygonal segments, and local obstacle surfaces.
  \item The far-field goal controller is structurally inspired by the ROS code, but it is not yet a line-by-line reproduction of the SO(3)-style geometric block.
  \item The magnetic term is mathematically faithful in structure, but the current implementation uses a compact 2D analogue rather than a direct symbolic port of the original 3D cross-product expressions.
  \item The cleanroom model removes all ROS-specific concerns: TF, publishers/subscribers, launch files, point-cloud nodes, timestamps, and platform adapters.
\end{enumerate}

\subsection*{Why This Change Was Made}

The goal of this first cleanroom step is verification by isolation:

\begin{itemize}[leftmargin=*]
  \item keep the algorithm small enough to reason about
  \item remove ROS coupling
  \item make rollout behavior inspectable with CSV and plot outputs
  \item establish a stable place for future refactors and comparisons
\end{itemize}

\subsection*{Required Follow-Up}

Future major changes should update this file, especially if they modify:

\begin{itemize}[leftmargin=*]
  \item the exact goal-relaxation gain
  \item the magnetic obstacle formula
  \item the near-goal switching logic
  \item the state update model
  \item the obstacle abstraction
  \item the benchmark scenario assumptions
\end{itemize}

\end{document}
